% RESEARCH_DRAFT_SHA256: d8306c7a24e8ea5a0125340e6016e3a2aae6e2c7b610dcdf92378306c40f675f
% Layout and source-reference wrappers expanded for the web renderer.
\documentclass[11pt]{article}
\usepackage[margin=25mm]{geometry}
\usepackage{amsthm}
\usepackage{fontspec,unicode-math}
\setmainfont{Latin Modern Roman}
\setmathfont{Latin Modern Math}
\usepackage{amsmath,amssymb,mathtools,xurl,hyperref}
\hypersetup{colorlinks=true,urlcolor=blue,linkcolor=blue}
\setcounter{secnumdepth}{0}
\setlength{\parindent}{0pt}
\setlength{\parskip}{5pt}
\setlength{\emergencystretch}{3em}
\newcommand{\sref}[2]{\hyperlink{source:#1:#2}{[S#2]}}
\newcommand{\cataloguescope}[1]{\paragraph{Source record.}#1}
\newcommand{\floor}[1]{\left\lfloor#1\right\rfloor}
\newtheorem{proposition}{Proposition}
\newtheorem{lemma}[proposition]{Lemma}
\begin{document}
\paragraph{Research manuscript.} OpenAI \texttt{gpt-6-astra}, reasoning \texttt{ultra}; 11 October 2026.
\paragraph{Preserved source body.} The complete catalogue mathematical text and sources follow with layout and source links expanded, including source qualifications and its historical preparation notes. New research begins at the section headed Research attempt.

\section{Polynomial-Time Estimation in Binary Latent-Variable Models}\label{statement}

\subsection{Definitions and mathematical statement}

\paragraph{Shared notation.}Standard number systems are $\mathbb N,\mathbb Z,\mathbb Q,\mathbb R,\mathbb C$. Specialized notation and index ranges follow the cited source and the definitions below.


A latent-variable model generates observed variables from unobserved variables and parameters $W$. Identifiability means different parameters produce different observed distributions, modulo the source\textquotesingle s allowed symmetries.
\paragraph{Record scope and qualifications.}
The catalogue title is retained as metadata. Named certificates, contracts, bounds or frameworks in that title are not asserted as results of the original source extract.
The model, sample-access assumptions, parameter encoding, and general identifiability conditions require the source.

Source-record limitation: The binary latent-variable model, parameterization, and general identifiability conditions are absent.
\paragraph{Statement recorded in the supplied source.}
Under general identifiability conditions for binary latent-variable models, do statistically consistent procedures exist for estimating the parameter matrix $W$ in time polynomial in the number $d$ of hidden variables? \href{https://doi.org/10.4171/owr/2018/29}{[S2]}
\paragraph{Source status.}
The supplied catalogue labels this source record open; that metadata was checked against the supplied snapshot on 2026-10-11. The label is preserved for provenance and is not an independent verification of present open status or mathematical correctness. The catalogue literature assessment is dated 2026-08-21.
\subsection{Short English statement}
Estimate the binary latent-model parameter matrix consistently with polynomial dependence on the hidden dimension.
\subsection{Sources}
\begin{itemize}
\item[S1] Ulam AI and the UnsolvedMath Contributors, UnsolvedMath: A Curated Collection of Open Mathematics Problems, catalogue entry. Catalogue curation and metadata: CC BY 4.0. The source extract is retained; definitions, English summaries and qualifications are editorial additions.
\url{https://huggingface.co/datasets/ulamai/UnsolvedMath}.
\item[S2] Matrix Estimation Meets Statistical Network Analysis: Extracting low-dimensional structures in high dimension (2019), Oberwolfach Reports, DOI 10.4171/owr/2018/29, pp. 1745--1783
\url{https://doi.org/10.4171/owr/2018/29}.
\item[S3] Jaffe et al., ICML 2018
\url{https://proceedings.mlr.press/v80/jaffe18a.html}.
\item[S4] Jaffe et al., preprint
\url{https://arxiv.org/abs/1802.09656}.
\item Source-specific Oberwolfach Reports licences, including CC BY-SA 4.0 where applicable, are retained; see THIRD\_PARTY\_NOTICES.md.
\end{itemize}
\paragraph{Preparation and review.}
The original source and UnsolvedMath contributors retain attribution for the source extract. Codex (GPT-6; exact serving revision not exposed) assisted with definitions, English summaries, scope notes and formatting. No new human mathematical authorship or independently verified solution is asserted. Catalogue identity checked against the complete supplied registry and current repository. Distinct upstream records, including repeated questions and broad or damaged extracts, are retained for complete coverage of the source CS category. A repeated record is not a new independent conjecture. Complete mathematical-equivalence and current-literature review has not been performed. Preparation and final review follow contribution prompts C and D, with source gaps disclosed. Structural and rendering checks do not establish a new theorem or open status.
\label{end:statement}

\section{Research attempt}
\subsection{Literature, source context and scope}
The full original catalogue body and source qualifications are preserved above. The relevant primary workshop passages were read in the cached publisher report. Historical claims and subsequent attributed results are distinguished from the mathematical work below. A restricted example, standard lemma or known theorem is not counted as new progress on the entire original target. Missing assumptions and repeated or rhetorical extracts are not silently turned into different solved problems.

\subsection{Identifiability does not supply a uniform sample bound}
The primary report restores the model absent from the short extract: $x=W^Th+\sigma\xi$, with $h\in\{0,1\}^d$, $d\le m$ and $W$ of row rank $d$. Its tensor method has exponential dependence on $d$ because it enumerates tensor eigenvectors. We test whether identifiability alone can justify a sample bound polynomial in $d$.
Take the noiseless submodel $m=d\ge1$, $\sigma=0$, and independent hidden bits, with $\Pr(h_1=1)=p=2^{-d}$ and all other bits unbiased. Compare $W_0=I_d$ and $W_1=\operatorname{diag}(2,1,\ldots,1)$. Both have full rank and their observed distributions are distinct, hence the two-point parameter family is identifiable.
\begin{proposition}
The total variation distance between the distributions of $N$ independent observations under these two parameters is $1-(1-p)^N\le Np$. Any estimator that, under both parameters, estimates the first diagonal entry with absolute error less than $1/2$ with probability at least $2/3$ therefore needs $N\ge2^d/3$.
\end{proposition}
\begin{proof}
When every sampled first hidden bit is zero, the observation laws agree exactly; this event has probability $(1-p)^N$ in both worlds. On the complementary event, the first observed coordinate takes a nonzero value at least once. That value is 1 in the first world and 2 in the second, so the two complementary supports are disjoint. This decomposition gives the exact total variation distance. The bound by $Np$ follows from the union bound over the rare-bit events.
The success intervals for the two first diagonal entries are disjoint. Turn the estimator into a test by accepting the first world whenever its estimate lies within $1/2$ of 1. This event has probability at least $2/3$ in the first world and at most $1/3$ in the second, so total variation is at least $1/3$. Combining inequalities yields $Np\ge1/3$ and the stated sample lower bound.
\end{proof}
\subsection{Verification and remaining gaps}
For small $d,N$, every sequence of first-bit observations was enumerated to obtain $1-(1-2^{-d})^N$ exactly, and the bound was checked. This rules out an inference from bare identifiability to a uniform polynomial-in-$d$ sample guarantee. It does not disprove the source's computational target: an algorithm polynomial in the number of supplied samples and in $d$ could still be consistent for every fixed identifiable model as sample size tends to infinity. Quantitative nondegeneracy may also exclude this rare-component family. No such algorithm or computational lower bound is obtained.

\section{References and source attribution}
Primary workshop report: \url{https://doi.org/10.4171/owr/2018/29}. Relevant passages read: 2018/29 p. 1770, x=W^T h+sigma xi, full row rank, tensor eigenpair enumeration and the polynomial-d question.

\url{https://proceedings.mlr.press/v80/jaffe18a.html}. Read primary ICML2018 abstract on moment/tensor estimation, nondegeneracy and consistency; the cached report supplies the explicit model and exponential eigenpair-enumeration dependence.

All retrievals were checked on 11 October 2026. Original question attribution remains in the preserved Sources section. External theorems are attributed and not claimed as new discoveries or fully verified merely by citation.


\paragraph{Integration review.} Codex (GPT-6; exact serving revision not exposed) reviewed the written arguments and expanded layout and source-reference wrappers. Repository metadata and the production web renderer were checked using Python 3.13 and JavaScript with MathJax 3.2.2. This model-assisted review is not independent human proof verification; no PDF compilation or formal proof checking is asserted.

\section{Completion Estimate}
\noindent\textbf{COMPLETION: 0\%}
No substantiated new progress toward the entire original target is claimed. The worked results and failed route are delimited above; a known result, illustrative model, or auxiliary obstruction is not counted as original-target completion.

\section{Final status and attribution}
\textbf{Attempt status: unresolved.} This describes this research attempt, not a new ruling on the source catalogue's status or on external literature claims.
The mathematics and this self-audit were generated by OpenAI \texttt{gpt-6-astra}, reasoning \texttt{ultra}, on 11 October 2026. Codex (GPT-6; exact serving revision not exposed) orchestrated the batch, reviewed the written mathematical arguments, and contributed formatting/repository review. This additional model review is not independent human verification. No human mathematical authorship, independent human proof checking, or formal verification is claimed. The historical preparation paragraphs in the reproduced catalogue body concern that body only.

Contribution prompt A was applied to the complete definition and supplied source record, followed by the adversarial checks of prompt B. Local reproduction files are \texttt{batch022/checks.py}, \texttt{computations.json}, and the per-ID audit JSON. They were not public links at preparation time. Numerical checks are finite and do not prove unrestricted results. No LaTeX engine was found locally; compilation was not performed by this research agent.

\end{document}
